% Source: https://artofproblemsolving.com/community/c3h1844631_equations_with_a_higher_degree
% Stable identity: aops-topic-1844631-q03
Let $A=2x^2-1$, $B=x^2-3x-2$ and $d=2x+4$. Then
\[
2x^2+2x+3=A+d,\qquad x^2-x+2=B+d,
\]
so the equation reads $\sqrt A+\sqrt B=\sqrt{A+d}+\sqrt{B+d}$, where $A,B\ge0$ for the left side to be defined, and $A+d$, $B+d$ are always positive (their discriminants $4-24$ and $1-8$ are negative).

If $d>0$, then $\sqrt{A+d}>\sqrt A$ and $\sqrt{B+d}>\sqrt B$, so the right side is larger. If $d<0$, both inequalities reverse and the right side is smaller. So we need $d=0$, i.e.\ $x=-2$.

For $x=-2$: $A=7$, $B=4+6-2=8$, and both sides equal $\sqrt7+\sqrt8$. So the only solution is $\boxed{x=-2}$.
